(1) 取单位面积控制体，稳态连续方程给出质量通量 \(j=\rho_1u_1=\rho_2u_2>0\)，量纲为 \(\mathrm{kg\,m^{-2}\,s^{-1}}\)。令比容 \(v_i=1/\rho_i\)，则 \(u_i=jv_i\)。动量通量连续意味着
\[
 p_1+\rho_1u_1^2=p_2+\rho_2u_2^2,\qquad
 p_2-p_1=j^2(v_1-v_2).
\]
两侧无热、无轴功，总焓通量守恒；对热完全理想气体 \(h=c_pT=\gamma pv/(\gamma-1)\)，所以
\[
 h_2+\frac{u_2^2}{2}=h_1+\frac{u_1^2}{2},\qquad
 h_2-h_1=\frac{j^2}{2}(v_1^2-v_2^2)
 =\frac{p_2-p_1}{2}(v_1+v_2).
\]
最后一个式子是 Hugoniot 关系，上一式为 Rayleigh 关系。这里的“绝热”指控制体没有净热输入，不意味着激波内可逆；真实压缩激波内部的黏性和导热使熵增加，因此不能把 \(pv^\gamma=\mathrm{const}\) 当作跨激波守恒律。

(2) 用 \(u_1^2=\gamma(p_1/\rho_1)x\) 及 \(u_2=u_1/r\) 化为无量纲方程
\[
 P-1=\gamma x(1-r^{-1}),\qquad
 \frac{\gamma}{\gamma-1}\left(\frac{P}{r}-1\right)
 =\frac{\gamma x}{2}(1-r^{-2}).
\]
消去 \(P\)，除去无激波的平凡根 \(r=1\)，得压缩激波支
\[
 \boxed{r=\frac{(\gamma+1)x}{(\gamma-1)x+2}},\qquad
 \boxed{P=\frac{2\gamma x-(\gamma-1)}{\gamma+1}}.
\]
由 \(p=\rho RT\) 得
\[
 \boxed{\frac{T_2}{T_1}=\frac{P}{r}
 =\frac{[2\gamma x-(\gamma-1)][(\gamma-1)x+2]}{(\gamma+1)^2x}}.
\]
声速平方为 \(c_2^2=\gamma p_2/\rho_2=c_1^2P/r\)，所以
\[
 \boxed{M_2^2=\frac{u_2^2}{c_2^2}=\frac{x}{Pr}
 =\frac{(\gamma-1)x+2}{2\gamma x-(\gamma-1)}}.
\]
对 \(x>1\)，有 \(r-1=2(x-1)/[(\gamma-1)x+2]>0\)，\(P-1=2\gamma(x-1)/(\gamma+1)>0\)。同时 \(M_2^2\) 分母减分子为 \((\gamma+1)(x-1)>0\)，故 \(0<M_2<1\)。上游超声速流经压缩激波后成为亚声速流。

(3) 每单位质量的理想气体熵满足 \(ds=c_v\,d\ln T+R\,d\ln v\)，其中 \(c_v=R/(\gamma-1)\)。故
\[
 \boxed{\frac{\Delta s}{R}=\frac{\ln P-\gamma\ln r}{\gamma-1}}.
\]
记 \(F(x)=\ln P-\gamma\ln r\)。直接代入上式求导并因式分解，
\[
 F'(x)=\frac{2\gamma(\gamma-1)(x-1)^2}
 {x[(\gamma-1)x+2][2\gamma x-(\gamma-1)]}>0\quad(x>1).
\]
因 \(F(1)=0\)，所以任意非零压缩激波均有 \(\Delta s>0\)，严格符合热力学第二定律。令 \(x=1+\varepsilon\)，在 \(x=1\) 处展开，前两阶系数均消失，得到
\[
 F(1+\varepsilon)=\frac{2\gamma(\gamma-1)}{3(\gamma+1)^2}\varepsilon^3+O(\varepsilon^4),\qquad
 \boxed{\frac{\Delta s}{R}=\frac{2\gamma}{3(\gamma+1)^2}\varepsilon^3+O(\varepsilon^4)}.
\]
因此弱激波的熵产生为激波强度的三次量，而非一次或二次量。

(4) 令 \(x=M_1^2\to\infty\)，可逐项取极限：
\[
 \boxed{r\to\frac{\gamma+1}{\gamma-1}},\qquad
 \boxed{\frac{P}{M_1^2}\to\frac{2\gamma}{\gamma+1}},\qquad
 \boxed{\frac{T_2/T_1}{M_1^2}\to\frac{2\gamma(\gamma-1)}{(\gamma+1)^2}},\qquad
 \boxed{M_2^2\to\frac{\gamma-1}{2\gamma}}.
\]
这说明理想气体强激波的密度压缩比有上界，而压力与温度仍按 \(M_1^2\) 增长。以 \(\gamma=7/5\)、\(M_1=2\) 检查，得 \(r=8/3\)、\(P=9/2\)、\(T_2/T_1=27/16\)、\(M_2^2=1/3\)；代入 (2) 两条无量纲守恒方程，左右严格相等。质量通量、压强和比焓的量纲分别为 \(\mathrm{kg\,m^{-2}\,s^{-1}}\)、\(\mathrm{Pa}\)、\(\mathrm{m^2\,s^{-2}}\)；上述比值与马赫数均无量纲。由于 \(u_2=u_1/r<u_1\)，跨激波比动能降低；无外部热功输入时，总焓守恒迫使比焓升高，内部不可逆耗散同时使熵增加。极限 \(M_1\downarrow1\) 时 \(r,P,T_2/T_1,M_2\to1\)，熵增趋零，与激波消失的物理边界一致。
